Cloud Gaming (CG) traffic generated by platforms such as GeForceNow, Microsoft Xbox Cloud Gaming, Sony PlayStation or Amazon Luna is foreseen to take an increasing and significant part of traffic shares in the upcoming years. However, CG traffic is particularly demanding because it requires at the same time a high downstream bandwidth to transmit the high-quality multimedia flow and a very low latency with no jitter to ensure a reactive response to the players' inputs. Previous studies \cite{marchal2023analysis} have shown that end-to-end mechanisms are not always sufficient to maintain a satisfactory QoS in some difficult network conditions. 

Meanwhile, the recent network technology of Active Queue Management (AQM) L4S (Low Latency Low Loss Scalable throughput) has been introduced and prioritizes the transport of low-latency (LL) traffic using a dedicated queue. However, identifying which flows should benefit from the LL queue remains a challenge. Addressing this issue, our previous work \cite{NOMS} proposed an initial set of VNF (Virtualized Network Functions) able to identify CG traffic at the network edge through a supervised Machine Learning (ML) model based on Decision Trees (DT) using statistical flow features. 

\contributions{
To summarize, the main contribution of this paper are as follows:
\begin{itemize}
    \item We propose a new model based on unsupervised ML that shows a far better capacity to recognize traffic from CG platforms that are missing in the learning dataset; 
    \item We translated the VNF facing the traffic to extract the features in P4 (Programming Protocol-independent Packet Processors) and deployed the module on a hardware Tofino based switch.
    \item In addition to those two contributions, we provide an extensive feedback on the adaptations needed to cope with the inherent limitations of P4 programming for hardware switches that enable its guaranteed high bitrate processing.
\end{itemize}
}


% capacity.
 
% The rest of this paper is organized as follows. Section \ref{chap_4:related_work} presents the related work on traffic detection and classification using P4.
% Then, Section \ref{chap_4:method} describes our dataset, our previous supervised ML model and the new unsupervised model we propose.
% Section \ref{chap_4:implement} presents our hybrid classification architecture mixing pure software network functions with a new P4 module for feature extraction at line rate, with a particular focus on the limitations imposed by the latter and the required adaptations. Finally, we evaluate the performance of our architecture in Section \ref{chap_4:eval} before concluding this study in Section \ref{chap_4:conclusion}.